\chapter{Introduction}\label{ch:intro}

\section{Context}

A conventional SAR cannot improve its azimuth resolution and its swath width at the same time: a fine azimuth resolution
needs a short antenna and therefore a wide Doppler bandwidth, which must be sampled with a high PRF, while a wide swath
needs a long interval between pulses. Multichannel systems break this trade-off by receiving each pulse with several
apertures displaced along track \cite{Krieger2004,Gebert2009}. Each channel is sampled at a PRF below the Doppler bandwidth
and is therefore aliased, but the set of channels samples the azimuth signal at the full rate, and the unambiguous
spectrum is recovered by an \emph{azimuth reconstruction}: a bank of filters that combines the channels so that the
aliased parts cancel.

Multistatic formations distribute the receivers over several platforms. They lift the limit set by the length of a
single antenna and offer flexibility, but the baselines between platforms are no longer small \cite{Sakar2020tgrs,Sakar2022thesis}.
In particular, a \emph{cross-track} baseline makes the bistatic range history of a scatterer depend on its height. The
reconstruction filters are derived for a reference geometry (here: the reference point of the scene at height
$h_0=0$); a scatterer at another height carries a residual phase that is different in each channel. The filters then
no longer match the data, the aliased parts no longer cancel and azimuth ambiguities appear in the image.

The same kind of height-dependent mismatch is known from airborne motion compensation, where the compensation is
computed for a reference height and leaves residual errors for the rest of the topography. Prats et al.\ proposed the
sub-aperture topography- and aperture-dependent (SATA) algorithm to correct it with an external DEM
\cite{Prats2005topo,Prats2005thesis}: the azimuth signal is split into short sub-apertures, each Doppler frequency of a
sub-aperture is associated with a ground position, and the phase residual of that position, read from the DEM, is removed.

\section{Objective}

The fellowship topic is to investigate whether SATA can correct the azimuth ambiguities caused by topography and
cross-track baselines in multistatic SAR systems. During the work the question was made precise as:
\begin{enumerate}
  \item How large are the ambiguities caused by unaccounted topography, as a function of the terrain, the baselines and
        the number of channels?
  \item How much of them can be removed using \emph{only the system parameters and a DEM}, with no knowledge of where
        the scatterers are?
\end{enumerate}
The second condition is strict on purpose. Several intermediate methods tested during the project reached the
reference level only because their correction maps were built from the positions of the simulated targets; such
methods are not applicable to real data and are not counted as solutions here (\cref{sec:sata-maps}).

\section{Contributions}

\begin{itemize}
  \item A numerical simulator of the multistatic acquisition and of the multichannel reconstruction in one dimension
        (azimuth) and in two dimensions (azimuth and range, 3-D scenes), with along- and cross-track baselines and arbitrary
        terrain (\cref{ch:model,ch:recon}).
  \item An analysis of SATA applied to each aliased channel before the reconstruction, showing when it is exact and
        predicting the residual it leaves on general terrain (\cref{ch:sata}).
  \item A joint reconstruction that places the DEM phase inside the inversion, first per sub-aperture and finally over
        the whole line, solved by regularised least squares with conjugate gradients (\cref{ch:joint}), and its 2-D
        extension (\cref{ch:joint2d}).
  \item Systematic sweeps over terrains, numbers of channels and baselines, and an audit of every method with a DEM only
        (\cref{ch:results}).
\end{itemize}

\section{Outline}

\Cref{ch:model} introduces the acquisition geometry and the topographic residual $\dC$. \Cref{ch:recon} reviews the
multichannel reconstruction and explains how topography produces ambiguities. \Cref{ch:sata} describes SATA and its
limit for aliased channels. \Cref{ch:joint} presents the joint reconstruction and \cref{ch:joint2d} its 2-D version.
\Cref{ch:results} collects the results and \cref{ch:conclusions} concludes. \Cref{app:software} maps the methods to
the code in the repository.
